%! AP Statistics — Unit 1, Lesson 1.7 — Student Packet Cover Sheet
\documentclass[10pt]{article}
\usepackage{apstats-article}
\usepackage{apstats-boxes}

%=============================================================================
\begin{document}

% ── Full-bleed navy banner ────────────────────────────────────────────────────
\begin{tikzpicture}[remember picture, overlay]
  \fill[navy] (current page.north west)
    rectangle ([yshift=-0.9in]current page.north east);
\end{tikzpicture}
\vspace{-0.8in}

\begin{center}
  {\color{white}\textbf{\LARGE AP Statistics}}\\[4pt]
  {\color{skymid}\large\textbf{Unit 1: Exploring One-Variable Data}}\\[6pt]
  {\color{white}\textbf{\large Lesson 1.7 \quad Summary Statistics for a Quantitative Variable}}
\end{center}

\vspace{0.22in}

\namedateperiod

\vspace{0.18in}

% ── LEARNING TARGETS ─────────────────────────────────────────────────────────
\begin{learningtargetbox}
\small
\begin{enumerate}[leftmargin=1.5em, itemsep=5pt]
  \item \ldots{} calculate the \textit{mean} and \textit{median} for a quantitative
        dataset and interpret each statistic in context.
  \item \ldots{} compute the \textit{range}, \textit{quartiles} ($Q_1$, $Q_3$),
        and \textit{interquartile range} (IQR) for a dataset.
  \item \ldots{} apply the \textbf{$1.5 \times \text{IQR}$ fence rule} to identify
        outliers, and explain why a value is or is not an outlier.
  \item \ldots{} choose the \textit{appropriate} measure of center and spread
        (median + IQR vs.\ mean + standard deviation) based on the shape of the
        distribution and the presence of outliers.
\end{enumerate}
\end{learningtargetbox}

\vspace{0.14in}

% ── TABLE OF CONTENTS ────────────────────────────────────────────────────────
\begin{tocbox}
\small
\renewcommand{\arraystretch}{1.5}
\begin{tabularx}{\linewidth}{@{} c l X r @{}}
\rowcolor{sky}
\textbf{\#} & \textbf{Component} & \textbf{Description} & \textbf{Score} \\
\midrule
1 & Warm-Up        & Interpret a dotplot; estimate center and spread         & \blank{1.2cm} \\
2 & Guided Notes   & Mean, median, IQR, standard deviation, outlier rule     & \blank{1.2cm} \\
3 & Group Activity & Calculate and compare statistics; identify outliers      & \blank{1.2cm} \\
4 & Exit Ticket    & Compute summary statistics; apply the outlier rule       & \blank{1.2cm} \\
5 & Homework       & Multi-part practice: calculate, interpret, AP-style      & \blank{1.2cm} \\
\midrule
  & \multicolumn{2}{r}{\textbf{Total}} & \blank{1.2cm} \\
\end{tabularx}
\end{tocbox}

\vspace{0.14in}

% ── KEEP IN MIND ─────────────────────────────────────────────────────────────
\begin{remindbox}
\small
The statistics in this lesson make the SOCS framework from Lesson 1.6 precise.
``Center is approximately 25 minutes'' becomes ``The median is 25 minutes.''

\vspace{6pt}
\begin{tabularx}{\linewidth}{@{} >{\centering\arraybackslash\columncolor{goldbg}}p{0.06\linewidth}
                                    X @{}}
\textbf{\large!} &
\textbf{Mean is pulled toward the tail; median resists.}\enspace
In a skewed distribution, the mean follows extreme values. The median stays near
the typical center because it depends only on position, not magnitude. \\[6pt]
\textbf{\large!} &
\textbf{Always pair the right center with the right spread.}\enspace
If you choose median (resistant), pair it with IQR (also resistant).
If you choose mean, pair it with standard deviation. Never mix. \\[6pt]
\textbf{\large!} &
\textbf{The $1.5 \times \text{IQR}$ rule is mechanical, not absolute.}\enspace
A value flagged as an outlier should be \textit{investigated}, not automatically removed.
Check for data entry errors first; if the value is real, report it and explain. \\[6pt]
\textbf{\large!} &
\textbf{Standard deviation is always $\geq 0$.}\enspace
$s_x = 0$ only when every data value is identical. The units of $s_x$ match the
units of the original data --- it is a typical distance from the mean. \\
\end{tabularx}
\end{remindbox}

\end{document}
